\section{Method}
\label{sec:method}

These bottlenecks stem from how current DLLMs implement parallel decoding, not from parallelism itself. We introduce \textbf{Strided DLLM}---a flexible diffusion language model that decodes at any stride from 1 (exact AR) to $N$ (parallel diffusion), while preserving the AR model's causal structure, reusing its serving infrastructure, and minimizing additional compute.
\yifan{Moved from §3 closing per Fan's comment.}

\paragraph{Unified verify-and-decode training.}
Decoding in diffusion language models faces a fundamental tension: generating multiple tokens simultaneously introduces errors from the token-independence assumption~\citep{nie2024scalingmdlm}, yet correcting these errors requires the model to evaluate its own outputs---a capability that standard DLLM training never provides. A natural temptation is to manage uncertainty through heuristics: LLaDA~\citep{llada2025} defers low-confidence tokens via remasking, SDAR~\citep{sdar2025} applies multiple denoising iterations per block, and LLaDA 2.1~\citep{llada21_2025} adds a token-to-token editing stream that retroactively corrects committed tokens. These strategies reduce errors in practice, but they share a common limitation: verification relies on confidence thresholds rather than distributional guarantees, leaving the quality--throughput tradeoff fundamentally unresolved.

The root cause lies in how these models are trained. In all cases, supervision is restricted to \texttt{[MASK]} positions---the indicator $\mathbf{1}[x_t^i = \texttt{M}]$ ensures no gradient flows through clean tokens~\citep{llada2025}. The model learns to \emph{decode} (predict tokens at masked positions) but never to \emph{verify} (evaluate candidate tokens from a clean prefix). Without a learned verify distribution $p$, distribution-exact acceptance criteria such as $\min(1, p/q)$~\citep{leviathan2023speculative}---which guarantee output correctness regardless of draft quality---cannot be applied. LLaDA 2.1~\citep{llada21_2025} partially addresses this through its mixture of mask-to-token and token-to-token objectives, training the model to identify and correct wrong tokens. However, the bidirectional attention prevents KV cache reuse, and the dual-threshold ($\tau_\text{mask}$, $\tau_\text{edit}$) decoding remains heuristic.

This limitation is compounded by two structural gaps relative to AR models. First, the \emph{understanding--generation mismatch} identified in Section~\ref{sec:mot_quality}: switching attention patterns between prefill and decoding degrades the quality of the verify signal even if one were available. Second, generating $N$ tokens in parallel is inherently harder than sequential prediction---the independence assumption means the model cannot capture dependencies \emph{among} simultaneously generated tokens~\citep{nie2024scalingmdlm}, and errors compound as the model conditions on its own imperfect outputs~\citep{llada21_2025}.

AR models resolve both issues elegantly. Causal attention with next-token prediction trains a single distribution $p_\theta(x_{i+1} \mid x_{\leq i})$ that serves as both the generation signal and the verification anchor---at every position, on every training example. This is what makes speculative decoding~\citep{leviathan2023speculative} possible: the acceptance rate $\alpha = 1 - \mathrm{TV}(p, q)$ depends only on draft quality, while correctness is guaranteed unconditionally.

To restore this unification in a diffusion model, we observe that clean-token supervision \emph{requires} logit shift (otherwise position $i$ trivially observes token $i$), and logit shift \emph{requires} causal attention (bidirectional attention leaks future tokens). This chain---\emph{verify needs clean supervision $\to$ clean supervision needs logit shift $\to$ logit shift needs causal mask}---motivates our design. With causal attention and logit shift, every position (masked or clean) is trained with the same next-token cross-entropy loss (Eq.~\ref{eq:loss}): clean positions learn the causal anchor $p$, masked positions learn proposals $q$ aligned with $p$. We quantify this alignment with the \emph{introspective acceptance rate}---how often the model's own verify distribution would accept its decode distribution under the $\min(1, p/q)$ criterion (Table~\ref{tab:introspective_rate}).

\begin{table}[h]
\centering
\caption{\textbf{Introspective acceptance rate.} AR and our Strided DLLM achieve high self-agreement due to unified verify-and-decode training. Prior DLLMs trained without clean-token supervision show substantially lower rates.}
\label{tab:introspective_rate}
\small
\begin{tabular}{lccc}
\toprule
\textbf{Method} & \textbf{Attn} & \textbf{Verify Training} & \textbf{Accept Rate (\%)} \\
\midrule
Qwen3-8B (AR) & Causal & \ding{51} (inherent) & \textbf{XX} \\
Strided DLLM (Ours) & Causal & \ding{51} (logit shift) & \textbf{XX} \\
SDAR & Block-causal & \ding{55} & \textbf{XX} \\
LLaDA 2.0 & Bidir. & \ding{55} & \textbf{XX} \\
LLaDA 2.1 & Bidir. & Partial$^\dagger$ & \textbf{XX} \\
\bottomrule
\end{tabular}
\\[4pt]
{\scriptsize $^\dagger$ Mixture of M2T and T2T objectives; verification uses confidence thresholds, not $p/q$ acceptance.}
\end{table}

\subsection{Strided DLLM Training}
\label{sec:training}

We convert a pretrained AR model into a strided diffusion model that retains both the verify and decode capabilities identified above. The training recipe combines causal attention, logit shift, and an all-masked objective---each addressing one link in the chain from verification to generation.

\fan{try to merge some paragraph titles. Otherwise, it reads like a combination of many tricks. e.g., ``Causal attention with full masks'' and ``Training with logit shift''.}
\yifan{Merged into 2 paragraphs: causal training with logit shift + all-masked training objective. Revised intro to connect with verify-and-decode argument.}

\paragraph{Causal training with logit shift.}
We apply \emph{strict causal masking} uniformly across both the masked and clean portions: for any query position $j$ and key position $i$, attention is permitted only when $i \leq j$. Unlike SDAR~\citep{sdar2025} which uses block-causal attention (bidirectional within blocks), this holds for masked tokens in $x_t$ as well---each \texttt{[MASK]} attends only to preceding positions. In the clean region $x_0$, standard causal self-attention is used, identical to AR pretraining, preserving the understanding--generation unification (Section~\ref{sec:mot_quality}) and ensuring KV cache reuse at inference. We pair this with a logit shift: the hidden state at position $i$ is trained to predict token $i{+}1$ rather than token $i$. Standard masked diffusion trains $\text{hidden}[\texttt{M}_i] \to \text{token}[i]$, which breaks the AR model's inherent $\text{logits}[i] \to \text{token}[i{+}1]$ mapping and---crucially---prevents clean positions from providing a meaningful verify signal. Our shifted formulation preserves this mapping: clean positions produce the causal anchor $p$ (the verify distribution), while masked positions produce proposals $q$ (the decode distribution), both trained with the same next-token objective. We visualize the resulting attention mask in Appendix~\ref{app:attn_mask}.

\paragraph{All-masked training objective.}
Given a clean sequence $x_0 = (x_1, \ldots, x_L)$, we replace \emph{all} tokens with a special \texttt{[MASK]} token to obtain the fully masked input $x_t = (\texttt{M}, \ldots, \texttt{M})$. The training input is the concatenation $[x_t \,|\, x_0]$, where $x_t$ is the all-masked sequence and $x_0$ provides the clean reference. Unlike standard masked diffusion training which masks a random fraction $r$ of tokens---wasting $(1{-}r)$ of the compute on unsupervised positions---our all-masked regime ensures every position contributes a useful training signal, eliminating the supervision dilution identified in Section~\ref{sec:mot_quality}. We use a cross-entropy loss with shifted labels on all positions:
\begin{equation}
    \mathcal{L} = -\frac{1}{|\mathcal{S}|}\sum_{\ell \in \mathcal{S}} \log p_\theta(x_0^{\ell+1} \mid [x_t, x_0]_{\leq \ell}),
    \label{eq:loss}
\end{equation}
where $\mathcal{S}$ is the set of non-padding positions across both the masked region $x_t$ and the clean region $x_0$, and $x_0^{\ell+1}$ is the shifted target (next token). This is the same cross-entropy objective used in AR pretraining, applied uniformly to all positions---no separate distillation loss or teacher model is required. On clean positions, the loss trains the \emph{verify} pathway: the model learns to produce the causal anchor distribution $p_\theta(x_{i+1} \mid x_{\leq i})$, recovering the exact AR training objective. On masked positions, the loss trains the \emph{decode} pathway: the model learns to produce proposals $q$ from \texttt{[MASK]} hidden states, enabling strided generation at stride $N > 1$. Because both pathways share the same objective, the model naturally aligns $q$ with $p$, maximizing the introspective acceptance rate (Table~\ref{tab:introspective_rate}). At inference time, the stride can be controlled adaptively per request, falling back to exact AR decoding (stride 1) whenever quality is prioritized over throughput.


\subsection{Introspective Strided Decoding (ISD)}
\label{sec:inference}

\begin{figure*}[t]
\begin{center}
\includegraphics[width=0.99\textwidth]{figures/comparison_v23.pdf}
\end{center}
\caption{\textbf{Comparison of decoding paradigms.} \textbf{Top:} AR decoding generates one token per step (stride$=$1) with causal attention and well-optimized infrastructure. \textbf{Middle:} Block Diffusion uses bidirectional attention within blocks (stride$=N$), requiring multiple denoising iterations and custom infrastructure. \textbf{Bottom:} Our Strided DLLM uses strict causal attention with adaptive stride ($1 < \text{stride} < N$), producing a quality-guaranteed token plus verified proposals per step. ISD introspects via the causal anchor $\pi_{\text{ar}}$; Residual ISD (R-ISD) additionally gates a LoRA residual $\pi_{\text{L}}$ for bit-for-bit lossless output. Our method is a drop-in replacement within AR serving infrastructure.}
\label{fig:inference_block_n}
\end{figure*}

Introspective Strided Decoding (ISD) is the inference counterpart to strided training. While the training recipe equips the model with the ability to produce meaningful predictions from \texttt{[MASK]} positions, ISD provides the mechanism to \emph{adaptively} select the effective stride at each generation step---balancing throughput and output quality without relying on heuristic confidence thresholds.

\fan{merge paragraph (titles) so that it is clean and shows some principled methods instead of a collection of tricks.}

\paragraph{Stride step.}
To decode at stride $N$, we append $N$ \texttt{[MASK]} tokens to the current prefix and run a single forward pass. Due to logit shift, this produces $N$ token predictions: the logits at the last clean position yield a \emph{quality-guaranteed} token (identical to an AR prediction), while the logits at each \texttt{[MASK]} position yield \emph{strided proposals} sampled from proposal distributions $q_k$. In a single forward pass, the model thus produces one exact token and $N{-}1$ proposals.

\paragraph{Introspection step.}
To verify the proposals, we fill them into the sequence (replacing the \texttt{[MASK]} tokens) and run a second forward pass. Because the filled-in tokens are now clean, the logits at each position produce the \emph{causal anchor} distributions $p_k$---the true AR next-token distributions conditioned on the actual token sequence. This introspection step simultaneously generates the next $N$ \texttt{[MASK]} proposals for the following stride, amortizing the verification cost.

\paragraph{Adaptive stride via $p/q$ acceptance.}
Inspired by the acceptance criterion in speculative decoding~\citep{leviathan2023speculative}, we compare each strided proposal against its causal anchor: proposal $x_k$ sampled from $q_k$ is accepted with probability $\min(1, \, p_k(x_k) / q_k(x_k))$. When accepted, the proposal provably follows the causal anchor distribution. On rejection, the token is resampled from the corrected distribution $\text{normalize}(\max(0, p_k - q_k))$ and all subsequent proposals are discarded. This mechanism naturally adapts the effective stride: in high-confidence regions the model accepts all $N$ proposals (achieving stride $N{+}1$ with a bonus token), while in uncertain regions it falls back to shorter strides. The formal procedure is given in Algorithm~\ref{alg:spec_verify}.

\begin{algorithm}[t]
\caption{Introspective Strided Decoding (one iteration)}
\label{alg:spec_verify}
\begin{algorithmic}[1]
\REQUIRE Prefix tokens $x_{1:L}$, stride $N$, model $M$ with logit shift
\ENSURE Accepted new tokens appended to prefix
\STATE \textcolor{magenta}{// Stride step: one forward pass with MASK tokens}
\STATE Construct input $[x_{1:L},\, \texttt{MASK},\, \ldots,\, \texttt{MASK}]$ \hfill ($N$ masks)
\STATE Run $M$ $\rightarrow$ $\text{logits}_{\text{stride}}[1{:}L{+}N]$
\STATE Sample $x_{L+1} \sim \text{softmax}(\text{logits}_{\text{stride}}[L])$ \hfill \COMMENT{quality-guaranteed (AR prediction)}
\FOR{$k = 2$ \TO $N$}
    \STATE Sample $x_{L+k} \sim q_k \triangleq \text{softmax}(\text{logits}_{\text{stride}}[L{+}k{-}1])$ \hfill \COMMENT{strided proposal}
\ENDFOR
\STATE
\STATE \textcolor{magenta}{// Introspection step: one forward pass with filled-in tokens}
\STATE Construct input $[x_{1:L},\, x_{L+1},\, \ldots,\, x_{L+N}]$
\STATE Run $M$ $\rightarrow$ $\text{logits}_{\text{anchor}}[1{:}L{+}N]$
\STATE $p_k \leftarrow \text{softmax}(\text{logits}_{\text{anchor}}[L{+}k{-}1])$ for $k = 1, \ldots, N$ \hfill \COMMENT{causal anchor distributions}
\STATE
\STATE \textcolor{magenta}{// Introspective verification with adaptive stride}
\STATE $n_{\text{accepted}} \leftarrow 1$ \hfill \COMMENT{$x_{L+1}$ always accepted (quality-guaranteed)}
\FOR{$k = 2$ \TO $N$}
    \STATE Draw $r \sim \text{Uniform}(0, 1)$
    \IF{$r < \min\!\big(1,\; p_k(x_{L+k}) \,/\, q_k(x_{L+k})\big)$}
        \STATE Accept $x_{L+k}$;\quad $n_{\text{accepted}} \leftarrow n_{\text{accepted}} + 1$
    \ELSE
        \STATE Resample $x_{L+k} \sim \text{normalize}\!\big(\max(0,\; p_k - q_k)\big)$
        \STATE $n_{\text{accepted}} \leftarrow n_{\text{accepted}} + 1$;\quad \textbf{break}
    \ENDIF
\ENDFOR
\STATE
\STATE \textcolor{magenta}{// Bonus token if all proposals accepted}
\IF{$n_{\text{accepted}} = N$}
    \STATE Sample $x_{L+N+1} \sim \text{softmax}(\text{logits}_{\text{anchor}}[L{+}N])$
    \STATE $n_{\text{accepted}} \leftarrow n_{\text{accepted}} + 1$
\ENDIF
\STATE
\STATE Append $x_{L+1:L+n_{\text{accepted}}}$ to prefix; commit KV cache
\end{algorithmic}
\end{algorithm}

\paragraph{Lossless ISD with residual adaptation.}
The default strided DLLM uses full fine-tuning, which modifies all model weights. While ISD guarantees that the output follows the converted model's distribution, this distribution inevitably differs from the original base AR model. Training with \emph{LoRA} instead creates a new opportunity: the base weights remain untouched, and the LoRA residual can be selectively activated per token.

Inspired by the gated LoRA approach in multi-token prediction~\citep{applemtp2025}, we gate the LoRA residual with a per-token binary mask: \texttt{[MASK]} (proposal) positions use base+LoRA weights to produce high-quality strided proposals, while verify (introspection) positions use \emph{base-model-only} weights. Critically, because the entire model uses strict causal attention, verify positions cannot attend to proposal positions---the causal anchor distribution $p$ is computed from base-only weights over a base-only KV cache, identical to a pure base AR forward pass. This guarantees that ISD verifies against the \emph{exact} base AR distribution, making the output bit-for-bit lossless.

\paragraph{Distinction from speculative decoding.}
While ISD draws inspiration from the $p/q$ acceptance criterion, it differs from speculative decoding in important ways. First, ISD uses a \emph{single model} as both proposer and verifier---there is no separate draft model to train, maintain, or synchronize. Second, the stride capability is a \emph{training-time property}: the model is explicitly trained to produce high-quality proposals from \texttt{[MASK]} positions, rather than relying on a weaker model's approximations. Third, the effective stride adapts naturally per step based on the model's own confidence in its proposals, providing a smooth quality--throughput trade-off without manual tuning.

\paragraph{Theoretical speedup analysis.}
Let $p_k$ denote the acceptance probability of the $k$-th strided proposal, and $P_k = \prod_{j=1}^{k} p_j$ the cumulative acceptance probability. For ISD at stride $N$, the expected tokens per forward pass (TPF) is (see Appendix~\ref{app:tpf_derivation} for derivation):
\begin{equation}
    \text{TPF}_N = \frac{2 + P_1 + P_2 + \cdots + P_{N-2}}{2 - P_{N-1}}.
    \label{eq:tpf}
\end{equation}
At perfect acceptance ($p_k = 1$), $\text{TPF}_N = N$, recovering the theoretical maximum. At $p_k = 0$, $\text{TPF}_N = 1$, degenerating to AR. In the memory-bound decode regime, forward pass latency is approximately constant regardless of stride size, so wall-clock speedup $\approx$ TPF. With typical acceptance rates of $p \geq 0.85$, stride $N{=}3$ achieves $\text{TPF} \approx 2.3\text{--}2.4\times$ with a compute overhead of only ${\sim}2\times$---meaning ISD translates most of the parallel decoding benefit into real throughput gain with minimal wasted compute. \fan{``our evaluations show that FRAMEWORK name achieves close-to-optimal speedup empirically.''}

\subsection{Infrastructure: AR-Compatible Serving}
\label{sec:infra}

As identified in Section~\ref{sec:mot_systems}, prior DLLMs require custom inference pipelines that forgo the optimizations accumulated in AR serving stacks. Because Strided DLLM preserves strict causal attention, we integrate directly into SGLang~\citep{zheng2024sglang} as a drop-in extension---each ISD step maps to SGLang's native \emph{extend} mode (appending $2N{-}1$ tokens with causal attention) rather than requiring a custom denoising loop. The model immediately inherits paged KV cache management, continuous batching, and tensor parallelism without modification.

\paragraph{ISD step lifecycle.}
Each ISD decode step executes a tight loop: (1)~\emph{classify} each request as cold-start or verify based on the previous step's accept/reject state; (2)~fill \texttt{input\_ids} with the appropriate mix of committed, speculative, and \texttt{[MASK]} tokens; (3)~run the GPU forward pass (CUDA graph replay); (4)~\emph{verify} proposals against anchor logits using the fused Triton kernel; (5)~\emph{trim} KV cache slots for rejected and MASK positions; (6)~assemble output tokens and state for the next step. Steps 4--6 require the GPU forward's output, creating a strict dependency chain: $\text{forward}(t) \to \text{verify}(t) \to \text{trim}(t) \to \text{prepare}(t{+}1) \to \text{forward}(t{+}1)$. Unlike AR, where CPU scheduling overlaps with GPU compute on a separate CUDA stream, this chain forces all CPU work onto the critical path.

\paragraph{Extend mode and CUDA graph capture.}
Extend mode is fundamentally heavier than AR's decode mode: attention metadata (page tables, KV index pointers, flashinfer plans) must be rebuilt each step (${\sim}$0.43\,ms vs.\ 0.07\,ms for decode). To eliminate the dominant cost---kernel launch overhead at small batch sizes---we capture the entire DLLM extend forward (all $2N{-}1$ tokens, including attention and LM head) into a single CUDA graph. On replay, only \texttt{input\_ids} and attention metadata are updated; all other kernel launches, memory allocations, and stream operations are replayed from the captured trace. This reduces forward time from ${\sim}$14.7\,ms (eager) to ${\sim}$7.1\,ms (graph replay) at $C{=}1$, the single largest latency reduction. For extend mode, we use \emph{paged-only attention} (a single flashinfer kernel per layer with causal masking) rather than the cascade of ragged + paged + merge kernels used in standard extend, eliminating 72 kernel launches per step (2 extra per layer $\times$ 36 layers).

\paragraph{KV cache management for speculative positions.}
Unlike AR decoding (which only appends KV entries), ISD must \emph{free} KV slots after each verification step: rejected proposals and trailing MASK positions no longer hold valid hidden states. After verify, the algorithm computes a \emph{trim count} (positions to free from the end of the KV segment) and an \emph{advance} (positions to commit). The output processor frees the trimmed KV indices in a single batched operation and advances each request's committed length. This trim-and-commit cycle is unique to speculative strided decoding and runs on every step, so we batch the free operations across requests to avoid per-request GPU synchronization.

We introduce four further optimizations that address the unique overhead structure of strided decoding; their cumulative effect is quantified in the ablation (Table~\ref{tab:ablation}).

\paragraph{Stationary-batch decode loop.}
SGLang's standard scheduling pipeline rebuilds the batch object each step: caching unfinished requests (a GPU tensor copy per request), filtering and merging batches, and constructing new attention metadata from scratch. For ISD, this runs on the critical path---unlike AR, where CPU scheduling overlaps with GPU compute. At batch size 32, the per-step CPU overhead exceeds 4ms (64 GPU tensor copies plus Python batch reconstruction).

We bypass this with a \emph{stationary-batch inner loop} that reuses the batch object across consecutive ISD steps. The lightweight \texttt{prepare\_for\_dllm\_decode} allocates KV slots via a single batched scatter, updates sequence lengths arithmetically ($\text{prev} + B \times \text{block\_size}$), and reuses cached metadata: request pool indices and position offset tensors are constant when the batch composition is stable, and \texttt{input\_ids} buffers are pre-allocated and filled with \texttt{[MASK]} (the algorithm overwrites specific positions). Non-critical work---notably \texttt{stream\_output} to the detokenizer (${\sim}$120\,$\mu$s)---is deferred to an overlap window during the next step's GPU forward. The loop exits to the full scheduling path only when new requests arrive or the batch empties. This eliminates the $O(B)$ GPU tensor copies that dominated at high concurrency, yielding the single largest speedup: $+$71\% at $C{=}1$ and $+$134\% at $C{=}32$.

\paragraph{Acceptance-maximizing proposals.}
In speculative decoding, draft quality directly affects output quality because draft tokens may be emitted without verification. In ISD, every emitted token is either verified against the causal anchor $p$ or resampled from the correction distribution $\max(0, p - q)$---output quality is guaranteed regardless of proposal quality. This decouples proposal strategy from output correctness, creating a degree of freedom unique to ISD: proposals should maximize \emph{acceptance probability}, not diversity.

We exploit this by replacing stochastic sampling with argmax at proposal positions: $x_k = \arg\max q_k$ for $k \geq 2$, while the quality-guaranteed token at position 1 is still sampled with the user's temperature and top-$p$ settings. Because $p$ and $q$ are aligned through shared training (Section~\ref{sec:training}), the most probable token under $q$ is also most likely to be accepted under $p$. This increases the acceptance rate from 68\% to 78\% and improves TPF by ${\sim}$10\%, with no effect on output diversity: the acceptance criterion $\min(1, p/q)$ ensures that every accepted token provably follows the anchor distribution $p$, and rejected tokens are resampled from $\max(0, p - q)$ independently of the proposal method.

\paragraph{Fused introspective verification.}
The verification step---computing softmax over anchor logits, gathering $p(x_k)$ and $q(x_k)$, evaluating the $p/q$ acceptance ratio, and corrective resampling on rejection---would normally require 7+ separate GPU kernel launches per verified position. We fuse these into a single Triton kernel with a two-pass early-exit design. Pass~1 streams through the vocabulary with online softmax, computes $p(x_k)$, and evaluates acceptance in a single pass. On acceptance (78\% of positions), the kernel returns immediately without Pass~2. On rejection, Pass~2 streams through the vocabulary again, computing the correction distribution via Gumbel-max sampling: $\arg\max(\log\max(0, p - q) + G)$ where $G$ is Gumbel noise, producing the corrected token in one pass without explicit normalization or multinomial sampling. To avoid storing the full draft distribution (${\sim}$39MB at $C{=}32$), we retain only the scalar probability $q(x_k)$ per proposal, reducing memory from $O(BNV)$ to $O(BN)$.

\paragraph{Segment-gated LoRA for lossless ISD.}
For lossless ISD (Section~\ref{sec:inference}), different positions within a single forward pass require different model weights: verify positions must use base-only weights to produce the exact AR anchor $p$, while proposal positions use base+LoRA weights for high-quality drafts $q$. Standard LoRA applies the adapter uniformly to all positions, which would contaminate the verify distribution. We introduce a five-part optimization stack to make conditional LoRA practical at ISD speeds; Table~\ref{tab:lora_progression} shows the cumulative impact.

\emph{(i) Conditional mask.} We gate the LoRA residual with a per-token binary mask applied after the shrink projection: $\text{output} = Wx + (\text{mask} \odot Ax) B^\top \cdot \alpha$, where $\text{mask}[i] = 0$ for verify positions and 1 for proposal positions. The mask is a pre-allocated GPU tensor updated in-place each step by \texttt{\_setup\_conditional\_lora}, which reads the classify output (cold-start vs.\ verify, number of speculative tokens) to determine which positions are MASK. This enables CUDA graph capture---the \texttt{mul\_} operation is baked into the graph, and only the mask \emph{contents} change between replays.

\emph{(ii) CUDA graph capture with real LoRA weights.} Naively, CUDA graph captures LoRA kernels with empty adapter IDs (the default in SGLang), producing no-op kernels that do nothing on replay. We pre-load the LoRA adapter into GPU memory \emph{before} graph capture and pass real adapter IDs during capture, so cuBLAS operations are recorded with valid weight pointers. On replay, \texttt{prepare\_lora\_batch} is skipped entirely (the graph replays baked-in ops), and only the conditional mask is updated. This reduces forward time from 13.4\,ms (eager with 288 Triton csgmv kernels) to 8.3\,ms (graph replay).

\emph{(iii) cuBLAS replacement for small-$M$ LoRA.} At the token counts typical of ISD ($M = 2N{-}1 \leq 9$), SGLang's default Triton csgmv kernels are inefficient: they use 48 registers and 32\,KB shared memory per kernel, yielding 11--19\,$\mu$s per launch. We replace them with cuBLAS \texttt{torch.mm} (shrink) and \texttt{torch.addmm\_} (expand), which use 168 registers and 160\,KB shared memory, achieving 4.5--6.5\,$\mu$s per kernel---a 2--3$\times$ improvement. For multi-slice layers (QKV: 3 slices, GateUp: 2 slices), we split the expand into per-slice \texttt{addmm\_} calls with appropriate offset indexing. This reduces total LoRA overhead from 4.0\,ms to 1.3\,ms per forward.

\emph{(iv) Two-stream overlap.} The base projection $Wx$ and LoRA shrink $Ax$ read the same input $x$ but write to independent outputs, exposing data parallelism. We dispatch the shrink (plus mask application) to a dedicated CUDA stream, running it concurrently with the base matmul on the default stream:
{\footnotesize
\[
\underbrace{\texttt{mm}(x,A^\top) \to \texttt{mul\_}}_{\text{lora\_stream}} ~\|~ \underbrace{\texttt{mm}(x,W^\top)}_{\text{main}} ~\to~ \underbrace{\texttt{addmm\_}}_{\text{sync}}
\]
}
At $M{=}5$, the base projection uses $<$15\% of the H100's 132 SMs, leaving ample capacity for the concurrent shrink. Pre-allocated shrink output buffers (keyed by number of slices: 1, 2, or 3) avoid tensor allocation on the side stream during CUDA graph capture. This hides 57\% of LoRA shrink latency.

\emph{(v) Deferred draft softmax.} The standard verify path computes draft probabilities $q(x)$ via a full \texttt{F.softmax} before storing them for the next step's verification. We replace this with a fused path: the verify kernel takes raw logits as input and computes softmax on-the-fly during the acceptance check (Pass~1 already streams through the vocabulary). This eliminates two redundant \texttt{F.softmax} calls per step.

\begin{table}[t]
\centering
\caption{\textbf{Lossless ISD optimization progression} ($N{=}3$, $C{=}1$, 1$\times$H100). Each row adds one optimization cumulatively. The conditional mask reduces TPS (from 228 to 194) because verify positions now correctly use base-only weights, lowering the acceptance rate from 80\% to 71\%.}
\label{tab:lora_progression}
\small
\begin{tabular}{lcc}
\toprule
\textbf{Configuration} & \textbf{TPS} & \textbf{Mechanism} \\
\midrule
Baseline (csgmv in CUDA graph) & 177 & 288 Triton kernels, 11--19$\mu$s each \\
+ cuBLAS replacement & 228 & \texttt{mm}/\texttt{addmm\_}, 4.5--6.5$\mu$s each \\
+ Conditional mask & 194 & verify=base, MASK=LoRA (accept: 80\%$\to$71\%) \\
+ Two-stream overlap & 212 & Shrink hidden behind base matmul \\
+ Deferred softmax & 213 & Fused verify eliminates 2$\times$ \texttt{F.softmax} \\
\midrule
Non-LoRA reference & 267 & No adapter overhead \\
\bottomrule
\end{tabular}
\end{table}
